% Standalone LaTeX source. Compile twice with pdflatex.
% No external bibliography, figures, or input files are required.
% The finite arithmetic certificate is checked by verify.cpp and certificate.txt.
% All mathematical constants also appear in this source. The C++ verifier
% and its machine-readable certificate are supplementary, not TeX inputs.
\documentclass[11pt]{article}
\usepackage[T1]{fontenc}
\usepackage[utf8]{inputenc}
\usepackage[margin=1in]{geometry}
\usepackage{amsmath,amsthm,mathtools}
\usepackage{newtxtext,newtxmath}
\usepackage{microtype}
\usepackage{booktabs}
\usepackage{float}
\usepackage{needspace}
\usepackage{enumitem}
\usepackage[hidelinks]{hyperref}
\usepackage[capitalize,nameinlink,noabbrev]{cleveref}

\newtheorem{theorem}{Theorem}[section]
\newtheorem{lemma}[theorem]{Lemma}
\newtheorem{proposition}[theorem]{Proposition}
\theoremstyle{definition}
\newtheorem{definition}[theorem]{Definition}
\theoremstyle{remark}
\newtheorem{remark}[theorem]{Remark}
\crefname{lemma}{Lemma}{Lemmas}
\Crefname{lemma}{Lemma}{Lemmas}
\crefname{proposition}{Proposition}{Propositions}
\Crefname{proposition}{Proposition}{Propositions}

\newcommand{\R}{\mathbb R}
\newcommand{\FR}{\operatorname{FR}}
\newcommand{\pos}[1]{\left[#1\right]^+}
\newcommand{\ind}{\mathbf 1}
\newcommand{\Qpath}{\mathcal Q}
\newcommand{\Ppath}{\mathcal P}
\newcommand{\radset}{\mathcal R}
\newcommand{\paramset}{\mathcal A}
\newcommand{\certfile}{\texttt{certificate.json}}
\newcommand{\verifier}{\nolinkurl{verify.cpp}}
\DeclareMathOperator{\dom}{dom}

\setlength{\emergencystretch}{2em}
\setlist[enumerate]{topsep=5pt,itemsep=4pt,parsep=0pt}
\setlength{\parskip}{3pt}
\allowdisplaybreaks[1]
\hypersetup{
  pdftitle={A 1.84345 Upper Bound for a Bidding Factor-Revealing Program},
  pdfsubject={A uniform bound for a thirteen-component hinge mixture, certified by exact arithmetic},
  pdfauthor={}
}

\title{A $1.84345$ Upper Bound for a\\
Bidding Factor-Revealing Program}
\author{}
\date{}

\begin{document}
\maketitle
\vspace{-2.5em}

\begin{abstract}
We give an explicit bidding function for which the factor-revealing program
restated in this note has value at most $36869/20000=1.84345$ in every
finite dimension. The bidding function is a rational convex combination of
thirteen discounted hinges. The proof allows arbitrary nonincreasing
radius rows, zero reference distances, and unordered diagonal entries.
A distribution on clock-to-radius thresholds averages valid facility
constraints and reduces the bound to inequalities for individual radius
paths. Convexity reduces these path inequalities to finitely many rational
endpoint calculations; a separate argument treats the unbounded parameter
ranges. The complete analytic reduction and all certificate constants are
given here. The accompanying exact-arithmetic verifier checks $120$
cumulative-mass inequalities and $112{,}802$ endpoint cases. An explicit
feasible family gives a nearly matching lower bound of
$1.843449818540222\ldots$ for the same bidding function.
\end{abstract}

\section{The program and the upper bound}\label{sec:program}

For a real number $t$, write $[t]^+=\max\{t,0\}$. Let
\begin{equation}\label{eq:bid}
  \beta(x,y)=\frac1{10^8}\sum_{s=1}^{13}W_s[x-\gamma_s y]^+,
  \qquad x,y\ge0,
\end{equation}
where the rational multipliers $\gamma_s$ and integer weights $W_s$ are
specified in \cref{tab:bid}. They satisfy $\gamma_s\ge1$, $W_s>0$, and
$\sum_sW_s=10^8$. Set
\begin{equation}\label{eq:rho}
  \rho=\frac{36869}{20000}=1.84345.
\end{equation}
All displayed finite decimals in this note denote exact rational numbers.

\begin{table}[H]
\centering
\caption{The bidding function. Each coefficient of $[x-\gamma_s y]^+$ is
$W_s/10^8$.}\label{tab:bid}
\renewcommand{\arraystretch}{1.1}
\setlength{\tabcolsep}{12pt}
\begin{tabular}{rrr@{\qquad}rrr}
\toprule
$s$ & $\gamma_s$ & $W_s$ & $s$ & $\gamma_s$ & $W_s$\\
\midrule
1 & $137/125$ & $66{,}678{,}856$ & 8 & $9/5$ & $659{,}515$\\
2 & $29/20$ & $616{,}823$ & 9 & $37/20$ & $2{,}091{,}320$\\
3 & $3/2$ & $12{,}834{,}612$ & 10 & $19/10$ & $2{,}245{,}119$\\
4 & $29/18$ & $4{,}845{,}692$ & 11 & $39/20$ & $2{,}936{,}362$\\
5 & $33/20$ & $1{,}922{,}894$ & 12 & $2$ & $766{,}770$\\
6 & $17/10$ & $70{,}695$ & 13 & $41/20$ & $2{,}724{,}139$\\
7 & $7/4$ & $1{,}607{,}203$ &  &  & \\
\bottomrule
\end{tabular}
\end{table}

\newpage

The following is the factor-revealing program of
\cite[Definition~4.1]{source-program}, restated in full. Its definition is
entirely in terms of real variables. The words \emph{client},
\emph{reference distance}, and \emph{radius} will serve only to describe
their roles in the proof; no underlying algorithm is needed.

\begin{definition}[Factor-revealing program]\label[definition]{def:program}
Fix an integer $n\ge 1$. The variables are a scalar $f$, reference distances
$d_1,\ldots,d_n$, and radii
\[
  r_{j,i}\qquad(1\le j\le n,\ j\le i\le n+1).
\]
Let $\FR_n(\beta)$ be the supremum of the objective
\begin{equation}\label{eq:objective}
  \Phi_\beta(r,f)
  =\sum_{j=1}^n\left(
       \sum_{i=j}^n\beta(r_{j,i},r_{j,i+1})+r_{j,n+1}
    \right)-f
\end{equation}
subject to
\begin{subequations}\label{eq:constraints}
\begin{align}
  \sum_{j=1}^n d_j&=1,
    &&\text{(normalization)},\label{eq:normalization}\\
  r_{j,i+1}&\le r_{j,i},
    &&1\le j\le i\le n,\label{eq:row}\\
  r_{i,i}&\le r_{j,i}+d_i+d_j,
    &&1\le j<i\le n,\label{eq:triangle}\\
  r_{j,\ell}&\le r_{k,\ell}+d_j+d_k,
    &&1\le j,k<\ell\le n+1,\quad j\ne k,\label{eq:metric}\\
  \sum_{j<i}\beta(r_{j,i},d_j)
       +\sum_{j=i}^n\beta(r_{i,i},d_j)&\le f,
    &&1\le i\le n,\label{eq:facility}\\
  d_j,\ f,\ r_{j,i}&\ge 0,
    &&\text{for all indices in range}.\label{eq:nonnegative}
\end{align}
\end{subequations}
\end{definition}

Thus \eqref{eq:row} is row monotonicity, \eqref{eq:triangle} is the triangle
constraint, \eqref{eq:metric} comprises the metric inequalities between
previous rows, and \eqref{eq:facility} comprises the facility constraints.
In particular, the definition imposes no ordering on the diagonal entries
$r_{j,j}$ and permits $d_j=0$.

\begin{theorem}[Uniform upper bound]\label[theorem]{thm:main}
For every integer $n\ge 1$, every choice of variables satisfying
\eqref{eq:row}--\eqref{eq:nonnegative} satisfies
\begin{equation}\label{eq:homogeneous-bound}
  \sum_{j=1}^n\left(
       \sum_{i=j}^n\beta(r_{j,i},r_{j,i+1})+r_{j,n+1}
    \right)
  \le f+\frac{36869}{20000}\sum_{j=1}^n d_j.
\end{equation}
Consequently,
\begin{equation}\label{eq:main-bound}
  \sup_{n\ge 1}\FR_n(\beta)\le\frac{36869}{20000}=1.84345.
\end{equation}
\end{theorem}

The homogeneous inequality \eqref{eq:homogeneous-bound} does not require
\eqref{eq:normalization}. With normalization, it is exactly an upper bound on
\eqref{eq:objective}. We prove it for the full feasible region, rather than
for any restricted family of radius rows.

\paragraph{Proof structure.}
The analytic argument reduces the theorem to two properties of a rational
probability distribution on $121$ thresholds: a cumulative-distribution condition and a finite
family of path bounds. These reductions occupy
\cref{sec:prelim,sec:sampling,sec:zero,sec:positive,sec:paths,sec:endpoints}.
\Cref{sec:certificate} states the finite certificate and completes the proof;
its data and verification procedure are given in \cref{app:verification}.
The verifier uses exact directed rounding to check the finite inequalities; the preceding lemmas establish their implication
for every dimension and every real feasible solution.

\section{Preliminary reductions}\label{sec:prelim}

\begin{lemma}[Properties of the bid]\label[lemma]{lem:properties}
The function $\beta$ is nonnegative, continuous, positively homogeneous,
nondecreasing in its first argument, nonincreasing in its second argument,
and jointly convex. Moreover,
\begin{equation}\label{eq:budget}
  0\le\beta(x,y)\le[x-y]^+,\qquad
  \beta(x,0)=x,\qquad \beta(x,x)=0.
\end{equation}
For $x\ge y\ge z\ge0$, it satisfies
\begin{equation}\label{eq:superadditive}
  \beta(x,y)+\beta(y,z)\le\beta(x,z).
\end{equation}
\end{lemma}

\begin{proof}
Each hinge $b_\gamma(x,y)=[x-\gamma y]^+$, with $\gamma\ge1$, has the
claimed continuity, homogeneity, monotonicity, and convexity properties.
It also satisfies $0\le b_\gamma(x,y)\le[x-y]^+$ and
$b_\gamma(x,0)=x$. Since \eqref{eq:bid} is a convex combination of such
hinges, these statements imply \eqref{eq:budget} and the other listed
properties of $\beta$.

For completeness, each hinge also satisfies \eqref{eq:superadditive}.
If $b_\gamma(x,y)$ and $b_\gamma(y,z)$ are both positive, then their sum is
\[
  x-\gamma y+y-\gamma z
   =x-\gamma z-(\gamma-1)y
   \le[x-\gamma z]^+.
\]
If at most one is positive, monotonicity gives the same bound. Taking the
nonnegative combination in \eqref{eq:bid} proves
\eqref{eq:superadditive}.
\end{proof}

In particular, the bid has all the admissibility properties of
\cite[Definition~3.1]{source-program}. In the upper-bound argument below,
the breakpoints of the reference-bid function $r\mapsto\beta(r,1)$ are
exactly the thirteen values $\gamma_1,\ldots,\gamma_{13}$ in
\cref{tab:bid}. Some of these exceed $\rho$; their inclusion in the scalar
endpoint set will be important.

\subsection{Eliminating terminal radii}

For $x\ge s\ge0$, \eqref{eq:budget} gives
\begin{equation}\label{eq:terminal-replace}
  \beta(x,s)+s\le x=\beta(x,0).
\end{equation}
Replace every terminal radius $r_{j,n+1}$ by zero, leaving all other variables
unchanged. This preserves row monotonicity, makes the terminal metric
inequalities trivial, and changes no facility constraint. By
\eqref{eq:terminal-replace}, it cannot decrease the objective. Hence it
suffices to prove \cref{thm:main} under the assumption
$r_{j,n+1}=0$ for all $j$.

For the rest of the proof, write
\begin{equation}\label{eq:aj-qj}
  a_j=r_{j,j},\qquad
  q_j=\sum_{i=j}^n\beta(r_{j,i},r_{j,i+1}),\qquad r_{j,n+1}=0.
\end{equation}
Telescoping the inequality $\beta(x,y)\le x-y$ along a row yields
\begin{equation}\label{eq:q-a}
  q_j\le a_j.
\end{equation}

\subsection{Central radii and a monotone clock}

For $2\le i\le n$, define
\begin{equation}\label{eq:center}
  m_i=\min_{j<i}(r_{j,i}+d_j).
\end{equation}
Choose $m_1>0$ sufficiently large that $a_1\le m_1+d_1$ and, when $n\ge2$,
$m_1\ge m_2$. Set $m_{n+1}=0$.

\begin{lemma}[Central-radius inequalities]\label[lemma]{lem:center}
The sequence $m_1,\ldots,m_{n+1}$ is nonincreasing and satisfies
\begin{align}
  r_{j,i}&\le m_i+d_j &&(1\le j\le i\le n),\label{eq:center-upper}\\
  m_i&\le r_{j,i}+d_j &&(1\le j<i\le n+1).\label{eq:center-lower}
\end{align}
\end{lemma}

\begin{proof}
For $2\le i\le n$, choose $p<i$ attaining the minimum in
\eqref{eq:center}. When $j=i$, \eqref{eq:triangle} gives
$r_{i,i}\le r_{p,i}+d_p+d_i=m_i+d_i$.
When $j<i$ and $j\ne p$, \eqref{eq:metric} gives the same inequality for
$r_{j,i}$. For $j=p$, it is immediate. The case $i=1$ follows from the choice
of $m_1$.
Inequality \eqref{eq:center-lower} follows directly from the minimum, or
from $m_{n+1}=r_{j,n+1}=0$ at the terminal column.

For $2\le i<n$, the same minimizer $p$ gives
\[
  m_{i+1}\le r_{p,i+1}+d_p\le r_{p,i}+d_p=m_i.
\]
The first and last steps follow from the choices of $m_1$ and $m_{n+1}$.
\end{proof}

Define the running maximum of the diagonal values by
\begin{equation}\label{eq:clock}
  T_0=0,\qquad T_i=\max_{p\le i}a_p\quad(1\le i\le n).
\end{equation}
The point of using $T_i$, rather than $a_i$, is that the former is
nondecreasing even if the diagonal values are not ordered.

\begin{lemma}[Facility constraints for the running maximum]\label[lemma]{lem:clock}
For every $1\le i\le n$,
\begin{equation}\label{eq:strengthened-facility}
  \widetilde B_i
   :=\sum_{j<i}\beta(r_{j,i},d_j)
     +\sum_{j\ge i}\beta(T_i,d_j)
   \le f.
\end{equation}
\end{lemma}

\begin{proof}
Choose $h\le i$ such that $a_h=T_i$, and compare with the original facility
constraint \eqref{eq:facility} at index $h$.
For $j<h$, row monotonicity gives $r_{j,i}\le r_{j,h}$.
For $h\le j<i$, it gives
\[
  r_{j,i}\le a_j\le T_i=a_h.
\]
For $j\ge i$, the first argument in
\eqref{eq:strengthened-facility} is exactly $T_i=a_h$.
Thus each summand is at most its counterpart in the facility constraint
at $h$.
\end{proof}

We now introduce a \emph{virtual process}, which is simply a convenient
way to order valid inequalities. Before event $i$, the central radius is
$m_i$, the rows with $j<i$ are called inactive and have radii $r_{j,i}$, and
the rows with $j\ge i$ are called active. A virtual clock $t$ increases from
$T_{i-1}$ to $T_i$. Inactive rows bid $\beta(r_{j,i},d_j)$ to the reference
facility, and active rows bid $\beta(t,d_j)$.
At event $i$, row $i$ becomes inactive, each row $j\le i$ updates to
$r_{j,i+1}$, and the central radius changes to $m_{i+1}$. The clock stays
at $T_i$ during the event. We call event $j$ the starting event of
row $j$ and $a_j$ its initial value.

Every clock-growth state, including its final pre-event state, has total
reference bid at most $f$, by \cref{lem:clock} and monotonicity of $\beta$.
An event only decreases that total: previous inactive radii decrease, and
the newly inactive row has radius at most $a_i\le T_i$.
The process begins with clock zero and positive central radius. At its end,
all radii and the central radius are zero.

\section{Averaging facility constraints}\label{sec:sampling}

Fix the thresholds
\begin{equation}\label{eq:times}
  u_\ell=1+\frac{\ell}{96},\qquad 0\le\ell\le120.
\end{equation}
All sums over threshold indices range over $\{0,\ldots,120\}$ unless
indicated otherwise. Let $\mu_0,\ldots,\mu_{120}$ be nonnegative masses summing to one. Their
exact values will be specified in \cref{app:verification}.

In the virtual process, the ratio $t/m$ is nondecreasing: the numerator is
nondecreasing and the denominator is nonincreasing. When $m=0$, define this
ratio to be $+\infty$. For each $u_\ell$, sample the process at its first
crossing of that threshold, according to the following convention.
During clock growth, use the state at which $t/m=u_\ell$. If instead an event
changes the ratio from below $u_\ell$ to at least $u_\ell$, use the state
\emph{immediately before that event}. The initial ratio is zero and the final
ratio is infinite, so every threshold has a sample. Several thresholds may
use the same pre-event state.

For row $j$, let $x_j(u_\ell)$ denote the first argument of its reference bid
at this sample: the clock value while active, or the current radius while
inactive. Define its \emph{reference credit} by
\begin{equation}\label{eq:credit}
  C_j=\sum_{\ell=0}^{120}\mu_\ell\beta(x_j(u_\ell),d_j).
\end{equation}
All sampled facility inequalities are valid. Averaging them gives
\begin{equation}\label{eq:credit-budget}
  \sum_{j=1}^n C_j
   =\sum_{\ell=0}^{120}\mu_\ell
       \sum_{j=1}^n\beta(x_j(u_\ell),d_j)
   \le f.
\end{equation}

We will prove, for each row separately, that
\begin{equation}\label{eq:row-goal}
  q_j\le\rho d_j+C_j.
\end{equation}
Once \eqref{eq:row-goal} is established, summing it and using
\eqref{eq:credit-budget} proves the homogeneous bound.
Sampling \emph{before} a threshold-crossing event is essential: it retains
the reference credit available before the radius decrease whose payment is
included in $q_j$.

\section{The zero-distance path inequality}\label{sec:zero}

Define the strict cumulative distribution
\begin{equation}\label{eq:cdf}
  H(u)=\sum_{\ell:\,u_\ell<u}\mu_\ell,
  \qquad H(+\infty)=1,
\end{equation}
and, for $u>1$, define
\begin{equation}\label{eq:kappa}
  \kappa(u)=\frac{\beta(u,1)}{u-1}
   =\sum_{s=1}^{13}\frac{W_s}{10^8}
       \left[1-\frac{\gamma_s-1}{u-1}\right]^+.
\end{equation}
Every summand is nondecreasing. Hence $\kappa$ is nondecreasing,
and $0\le\kappa(u)\le1$.

The first requirement on the masses is
\begin{equation}\label{eq:cdf-checks}
  \sum_{k<\ell}\mu_k\ge\kappa(u_\ell)
  \qquad(1\le\ell\le120).
\end{equation}
These are finitely many rational inequalities. They imply
\begin{equation}\label{eq:cdf-dominance}
  H(u)\ge\kappa(u)\qquad(u>1).
\end{equation}
Indeed, if $u_{\ell-1}<u\le u_\ell$, then
$H(u)=\sum_{k<\ell}\mu_k\ge\kappa(u_\ell)\ge\kappa(u)$.
Above $u_{120}$, the bound follows from $H(u)=1$.

The next lemma is an abstract statement about radius paths. We will use it
both for zero reference distances and for the unbounded parameter ranges
that arise later.

\begin{lemma}[Normalized zero-distance paths]\label[lemma]{lem:zero-path}
Assume \eqref{eq:cdf-checks}. Let
$1=z_0\ge z_1\ge\cdots\ge z_p=0$ be any finite path. For each $\ell$, let
\[
  \tau_\ell=\min\{t\in\{0,\ldots,p-1\}:z_{t+1}<1/u_\ell\},
  \qquad p_\ell(z)=z_{\tau_\ell}.
\]
Thus $p_\ell(z)$ is the radius just before the first strict drop below
$1/u_\ell$. Then
\begin{equation}\label{eq:zero-path-bound}
  \Qpath(z):=\sum_{t=0}^{p-1}\beta(z_t,z_{t+1})
   \le\Ppath(z):=\sum_{\ell=0}^{120}\mu_\ell p_\ell(z).
\end{equation}
\end{lemma}

\begin{proof}
Every $\tau_\ell$ exists because the final radius is zero.
Writing each pre-drop radius as the sum of its subsequent decrements gives
\begin{align}
  \Ppath(z)
   &=\sum_{t=0}^{p-1}(z_t-z_{t+1})
          \sum_{\ell:\,\tau_\ell\le t}\mu_\ell\notag\\
   &=\sum_{t=0}^{p-1}(z_t-z_{t+1})H(1/z_{t+1}),
       \label{eq:zero-decrements}
\end{align}
where $1/0=+\infty$. The second equality uses
$\tau_\ell\le t$ if and only if $z_{t+1}<1/u_\ell$.
For $0<y<x\le1$, homogeneity and \eqref{eq:cdf-dominance} imply
\[
  \frac{\beta(x,y)}{x-y}
    =\kappa(x/y)
    \le\kappa(1/y)
    \le H(1/y).
\]
For $y=0$, the corresponding inequality is an equality because
$\beta(x,0)=x$ and $H(+\infty)=1$. A zero decrement contributes zero to both
sides. Applying these inequalities to \eqref{eq:zero-decrements} proves
\eqref{eq:zero-path-bound}.
\end{proof}

\begin{lemma}[Rows with zero reference distance]\label[lemma]{lem:zero-client}
Assume \eqref{eq:cdf-checks}. If $d_j=0$, then $q_j\le C_j$.
\end{lemma}

\begin{proof}
If $a_j=0$, the claim follows from \eqref{eq:q-a}. Otherwise, divide row $j$
by $a_j$ to obtain a path $z$ as in \cref{lem:zero-path}.
After row $j$ becomes inactive, the clock is at least $a_j$, and
\eqref{eq:center-lower} gives $m\le r$ for its current radius $r$.
A drop of the normalized radius strictly below $1/u_\ell$ therefore forces
$t/m>u_\ell$. The actual sample at $u_\ell$ occurs no later than that event.

If the sample occurs while row $j$ is inactive, row monotonicity and the
pre-event sampling convention give
$x_j(u_\ell)\ge a_jp_\ell(z)$.
If it occurs while the row is active, its clock is at least $a_j$.
To see this for a crossing before its starting event, note that the central radius
is at least $m_j\ge a_j$; for a center-drop crossing, the same is true of the
post-event central radius. Thus the sampled clock is at least
$u_\ell m_j\ge a_j$. At the starting event itself, the clock is
$T_j\ge a_j$ as well.
Consequently, $x_j(u_\ell)\ge a_jp_\ell(z)$ in every case.
Using $\beta(x,0)=x$, homogeneity, and \cref{lem:zero-path}, we obtain
\[
  C_j\ge a_j\Ppath(z)\ge a_j\Qpath(z)=q_j.
\]
\end{proof}

\section{Positive-distance rows}\label{sec:positive}

Fix a row $j$ with $d_j>0$, and normalize its initial value and starting
central radius by $d_j$:
\begin{equation}\label{eq:normalized-a-m}
  a=\frac{a_j}{d_j},\qquad M=\frac{m_j}{d_j}.
\end{equation}
We call $a$ the normalized activation parameter.
When $a\le\rho$, \eqref{eq:q-a} already proves \eqref{eq:row-goal}.
Hence assume $a>\rho$. By \eqref{eq:center-upper},
\begin{equation}\label{eq:m-lower}
  M\ge a-1>0.
\end{equation}

Let $h$ be the number of thresholds satisfying $u_\ell\le a/M$.
We cannot have $h=121$: since $u_{120}=9/4$, that case would imply
\[
  \frac94\le\frac aM\le\frac a{a-1},
  \quad\text{and therefore}\quad a\le\frac95<\rho.
\]
Thus $0\le h\le120$. Since $u_h>a/M$, the actual value of $M$ is at least
\begin{equation}\label{eq:Mh}
  M_h(a)=\max\left\{a-1,\frac a{u_h}\right\}.
\end{equation}
For $h\ge2$, the condition $u_{h-1}\le a/M$, together with
\eqref{eq:m-lower}, also gives
\begin{equation}\label{eq:h-upper}
  a\le U_h:=\frac{u_{h-1}}{u_{h-1}-1}.
\end{equation}
For $h=0,1$, write $U_h=+\infty$.

\subsection{Lower bounds on the sampled credit}

For the first $h$ thresholds, define
\begin{equation}\label{eq:active-credit}
  A_h(a)=\sum_{\ell<h}\mu_\ell\beta(M_h(a)u_\ell,1).
\end{equation}
The threshold $u_\ell$ with $\ell<h$ has been reached by the state
immediately before the starting event of row $j$, because
$T_j/m_j\ge a/M\ge u_\ell$.
The row is still active at its sample.
Before row $j$ starts, the center is at least $m_j$. If an earlier center drop
crosses the threshold, its post-event center is also at least $m_j$.
In either type of crossing, the sampled clock is at least $u_\ell m_j$.
Thus the normalized reference bid at this sample is at least
$\beta(M_h(a)u_\ell,1)$.

For the remaining thresholds, put
\begin{equation}\label{eq:threshold-radii}
  b_\ell(a)=\frac a{u_\ell}-1\qquad(0\le\ell\le120).
\end{equation}
Consider $\ell\ge h$ with $b_\ell(a)>0$, and let $x\longrightarrow y$ be the
first normalized row drop with $y<b_\ell(a)$.
Immediately after this event, the normalized central radius is at most
$y+1<a/u_\ell$, by \eqref{eq:center-lower}, while the normalized clock is
at least $a$. Hence the actual sample at $u_\ell$ occurs no later than this
event.

If that sample is taken while the row is inactive, its normalized radius is
at least $x$. If it is taken while the row is active strictly before
its starting event, its normalized clock is at least $u_\ell M>a$, since
$u_\ell>a/M$. If it is taken at the starting event, its normalized clock
is at least $a$ as well. In all cases the normalized reference bid is at
least $\beta(x,1)$. For $b_\ell(a)\le0$, we use the lower bound zero.

Define
\begin{equation}\label{eq:spent}
  S_h(r;a)=\sum_{\substack{\ell\ge h\\b_\ell(a)>r}}\mu_\ell
  \qquad(r\ge0).
\end{equation}
This is the mass of remaining thresholds whose strict radius crossings have
occurred by radius $r$. It is nonincreasing as a function of $r$.
For a normalized row path
\[
  a=x_0\ge x_1\ge\cdots\ge x_p=0,
\]
the preceding bounds give
\begin{equation}\label{eq:credit-lower}
  \frac{C_j}{d_j}\ge A_h(a)+
     \sum_{t=0}^{p-1}
        \bigl(S_h(x_{t+1};a)-S_h(x_t;a)\bigr)\beta(x_t,1).
\end{equation}
Here $S_h(a;a)=0$, since $b_\ell(a)<a$ for every $\ell$.

The strict comparison in \eqref{eq:spent} is important. A drop that lands
exactly at $b_\ell(a)$ is allowed to defer that threshold's charge until a
later drop. A later pre-drop radius is no larger, so this convention only
decreases the reference-credit lower bound.

\subsection{The scalar path inequality}

For a finite nonincreasing path from $a$ to zero, define its excess by
\begin{equation}\label{eq:excess}
\begin{split}
  E_h(a;x_0,\ldots,x_p)
   ={}&\sum_{t=0}^{p-1}\Bigl[
      \beta(x_t,x_{t+1})\\[-2pt]
      &\hspace{3em}
      -\bigl(S_h(x_{t+1};a)-S_h(x_t;a)\bigr)\beta(x_t,1)
     \Bigr]-A_h(a).
\end{split}
\end{equation}
By homogeneity and \eqref{eq:credit-lower},
\begin{equation}\label{eq:row-excess-reduction}
  \frac{q_j-C_j}{d_j}\le E_h(a;x_0,\ldots,x_p).
\end{equation}
It remains to show that \eqref{eq:excess} is at most $\rho$ for every such
path and every $a\in[\rho,U_h]$. Including $a=\rho$ and any boundary values
not realizable by the original definition of $h$ only enlarges the domain
to be checked.

\section{A finite dynamic program for all radius paths}\label{sec:paths}

For fixed $a\ge\rho$, form the finite set
\begin{equation}\label{eq:radius-set}
  \radset(a)=
    \left(\{0,a,\gamma_1,\ldots,\gamma_{13}\}
       \cup\{b_\ell(a):0\le\ell\le120\}\right)\cap[0,a].
\end{equation}
List its distinct elements as
$0=v_0<v_1<\cdots<v_K=a$.
There are at most $136$ elements. Define
\begin{align}
  V_0(h,a)&=0,\notag\\
  V_k(h,a)
   &=\max_{0\le p<k}\Bigl\{
       \beta(v_k,v_p)
       -\bigl(S_h(v_p;a)-S_h(v_k;a)\bigr)\beta(v_k,1)
       +V_p(h,a)
     \Bigr\},\label{eq:dp}\\
  D_h(a)&=V_K(h,a)-A_h(a).\label{eq:D}
\end{align}
The recurrence computes the largest excess of a decreasing path through the
specified radii. The next lemma explains why it bounds paths with arbitrary
real radii and arbitrary length.

\begin{lemma}[Finite radius reduction]\label[lemma]{lem:radius-reduction}
For every $h\in\{0,\ldots,120\}$, every $a\ge\rho$, and every finite
nonincreasing path $a=x_0\ge\cdots\ge x_p=0$,
\begin{equation}\label{eq:radius-reduction}
  E_h(a;x_0,\ldots,x_p)\le D_h(a).
\end{equation}
\end{lemma}

\begin{proof}
Fix a path length. Partition each variable radius using the endpoints in
$\radset(a)$, allowing singleton cells at the endpoints. On the interior of
each resulting product cell, every threshold indicator in $S_h$ is fixed,
and each reference bid $\beta(x_t,1)$ is affine. The transition payments
$\beta(x_t,x_{t+1})$ are jointly convex by \cref{lem:properties}.
Consequently, on this cell the excess is a convex function of the radius
variables, and it extends convexly to the cell's closure by keeping the
threshold indicators fixed.

Intersect that closure with the constraints
$a=x_0\ge x_1\ge\cdots\ge x_p=0$.
This is a compact polytope, and a convex function on it is bounded above by
its largest vertex value. At any vertex, every radius equals an interval
endpoint. Indeed, a maximal block of equal coordinates that touches none of
its interval endpoints could be perturbed both upward and downward, by a
sufficiently small amount, without violating any interval or order
constraint. Such a point would not be a vertex. Thus every vertex radius
belongs to $\radset(a)$.

We next compare the convex extension's convention at a vertex with the
strict convention in \eqref{eq:spent}. For each threshold separately, a
limiting path may land exactly on its threshold radius. The strict
convention can only postpone that threshold's charge to a later drop;
it cannot move it earlier. Reference bids are nondecreasing in the radius,
so postponement weakly decreases the amount subtracted. The excess under
the strict convention is therefore at least its value under the convex
extension. If the threshold radius is zero, discarding its charge has the
same favorable effect.

Finally, remove consecutive repeated radii. They contribute no transition
payment because $\beta(x,x)=0$, and their strict $S_h$ values agree. The
remaining path is a decreasing path through $\radset(a)$. The recurrence
\eqref{eq:dp} considers every such path, and therefore bounds its excess.
This argument applies to every original path length, proving the lemma.
\end{proof}

\begin{remark}[What the reduction does not assume]
The set $\radset(a)$ is not a numerical grid imposed on the original
program. Its endpoints are determined by the bid's breakpoints and the
sampling thresholds. The convexity argument shows that these endpoints
suffice for an upper bound, without restricting the number of radius
changes in the original variables.
\end{remark}

\section{A finite set of activation parameters}\label{sec:endpoints}

For each $h\in\{0,\ldots,120\}$, let
\begin{equation}\label{eq:Ih}
  I_h=[\rho,U_h],\qquad
  U_h=\begin{cases}
       +\infty,&h=0,1,\\[2pt]
       \dfrac{u_{h-1}}{u_{h-1}-1},&2\le h\le120.
      \end{cases}
\end{equation}
Empty intervals are discarded. Define $\paramset_h$ by retaining all distinct
finite numbers in $I_h$ from the following list:
\begin{equation}\label{eq:endpoints}
\begin{gathered}
  \rho,\qquad U_h\ \text{when finite},\qquad
  \frac{u_h}{u_h-1}\ \text{when }u_h>1,\qquad
  \gamma_s\quad(1\le s\le13),\\[3pt]
  u_\ell,\qquad u_\ell(1+\gamma_s)
     \qquad(0\le\ell\le120,\ 1\le s\le13),\\[3pt]
  1+\frac{\gamma_s}{u_\ell},\qquad
  \frac{\gamma_su_h}{u_\ell}
     \qquad(0\le\ell<h,\ 1\le s\le13).
\end{gathered}
\end{equation}
The first line includes the domain endpoints, the branch switch of
$M_h(a)$, and the values where the initial radius $a$ crosses a
reference-bid breakpoint. The second includes the values where a threshold
radius $b_\ell(a)$ crosses zero or a reference-bid breakpoint. The third
includes the active-reference-bid breakpoints on both branches of $M_h(a)$.
All these values are rational. No breakpoint of a positive transition
payment is needed: those payments will be treated by convexity.

\begin{lemma}[Finite parameter reduction]\label[lemma]{lem:parameter-reduction}
Assume \eqref{eq:cdf-checks}. For every nonempty $I_h$,
\begin{equation}\label{eq:parameter-reduction}
  D_h(a)\le\max_{a'\in\paramset_h}D_h(a')
  \qquad(a\in I_h).
\end{equation}
\end{lemma}

\begin{proof}
We first treat bounded intervals between successive elements of
$\paramset_h$, then the two possible unbounded intervals.

\paragraph{Bounded intervals.}
On an open interval containing none of the breakpoints in
\eqref{eq:endpoints}, the order of the symbolic radius states
\[
  0,\quad a,\quad \gamma_1,\ldots,\gamma_{13},\quad
  a/u_0-1,\ldots,a/u_{120}-1
\]
is fixed, as is the subset lying in $[0,a]$. Indeed, two distinct moving
threshold radii never cross for $a>0$, and each is strictly below $a$.
A fixed reference breakpoint enters $[0,a]$ only at $a=\gamma_s$;
a threshold radius crosses zero at $a=u_\ell$ or a reference breakpoint
at $a=u_\ell(1+\gamma_s)$. Each such value appears in
\eqref{eq:endpoints}.

Each symbolic state is affine in $a$. Every threshold mass $S_h$ at such a
state is constant on the interval. Every inactive reference bid is affine
there: its affine piece can change only at $a=\gamma_s$ for the initial
radius, or at $a/u_\ell-1=\gamma_s$ for a threshold radius.
Both types were included in \eqref{eq:endpoints}.
The same holds for $A_h(a)$, since the branch of $M_h(a)$ is fixed and its
reference bids can change piece only at
\[
  (a-1)u_\ell=\gamma_s
  \quad\text{or}\quad a u_\ell/u_h=\gamma_s,
  \qquad \ell<h.
\]

Fix a path through the ordered symbolic states. Its positive transition
terms are convex functions of $a$, being nonnegative sums of positive
parts of affine functions of $a$.
All its subtracted terms are affine. Hence its excess is convex in $a$.
There are finitely many such paths, so their maximum $D_h(a)$ is also
convex on the open interval. It is therefore bounded by its one-sided
endpoint values.

At an endpoint, take the limit of any symbolic path. States may coincide,
or a threshold radius may become zero. As in the proof of
\cref{lem:radius-reduction}, the strict crossing convention weakly
increases its excess relative to any one-sided limit, and repeated states
can be removed. Thus the value of $D_h$ at the endpoint dominates the
one-sided limit. This proves the assertion on every bounded interval.

\paragraph{Unbounded intervals.}
Only $h=0,1$ have an unbounded domain. Beyond the last finite breakpoint,
the symbolic state order and all reference-bid pieces remain fixed.
For a fixed symbolic path, write its radii as $x_0(a),\ldots,x_p(a)$ and
its excess as $F(a)$. The function $F$ is convex on this ray.
The normalized limits
\[
  z_t=\lim_{a\to\infty}\frac{x_t(a)}a
\]
exist and form a nonincreasing path from one to zero. Every limit belongs
to
\[
  \{0,1\}\cup\{1/u_\ell:0\le\ell\le120\}.
\]
Homogeneity and continuity give
\begin{equation}\label{eq:payment-limit}
  \lim_{a\to\infty}\frac{\beta(x_t(a),x_{t+1}(a))}{a}
     =\beta(z_t,z_{t+1}),\qquad
  \lim_{a\to\infty}\frac{\beta(x_t(a),1)}{a}=z_t.
\end{equation}

Consider a remaining threshold $u_\ell$ and its first strict crossing in the
limiting path $z$, as defined in \cref{lem:zero-path}.
If $z_{t+1}<1/u_\ell$, then, for all sufficiently large $a$,
\[
  x_{t+1}(a)<a/u_\ell-1=b_\ell(a).
\]
Thus the finite path's charge for this threshold occurs no later than the
limiting path's first strict crossing. Its normalized pre-drop radius is
therefore at least $p_\ell(z)$ in the limit. In particular, equality with a
threshold in the limiting path cannot invalidate this bound: strict
crossing can only defer the limiting charge.

For $h=0$, all masses are remaining thresholds, so
\eqref{eq:payment-limit} implies
\[
  \limsup_{a\to\infty}\frac{F(a)}a
     \le\Qpath(z)-\Ppath(z)\le0
\]
by \cref{lem:zero-path}.
For $h=1$, the active contribution obeys
\[
  \lim_{a\to\infty}\frac{A_1(a)}a=\mu_0,
\]
since $M_1(a)=a-1$ for all sufficiently large $a$ and $u_0=1$.
This is precisely the missing zero-distance credit
$\mu_0p_0(z)$, because $p_0(z)=1$.
The same nonpositive bound on the asymptotic slope follows.

A convex function $F$ on a ray with
$\limsup_{a\to\infty}F(a)/a\le0$ is nonincreasing.
Otherwise it would have a positive secant slope, and convexity would force
its asymptotic slope to be at least that positive number.
Therefore each symbolic path excess is nonincreasing on the ray, and so
is their finite maximum. The finite left endpoint dominates the one-sided
limit by the bounded-interval argument. It consequently controls the
entire ray, completing the proof.
\end{proof}

The preceding two lemmas reduce all finite radius paths and every real
$a\in I_h$ to the values $D_h(a)$ at the finite set
\begin{equation}\label{eq:cases}
  \mathcal C=\{(h,a):0\le h\le120,\ a\in\paramset_h\}.
\end{equation}
No cutoff on large values of $a$ is used.

\section{The finite certificate and completion of the proof}\label{sec:certificate}

Write the sampling masses as
\begin{equation}\label{eq:mu}
  \mu_\ell=\frac{c_\ell}{10^9}\qquad(0\le\ell\le120),
\end{equation}
where the integers $c_\ell$ are given in \cref{tab:masses} in
\cref{app:verification}. They are nonnegative and sum to $10^9$.
These masses, the bid in \cref{tab:bid}, and the thresholds
\eqref{eq:times} constitute the rational certificate.

\begin{proposition}[Exact finite certificate]\label[proposition]{prop:certificate}
For the masses \eqref{eq:mu}, all $120$ inequalities
\eqref{eq:cdf-checks} hold. The set $\mathcal C$ in \eqref{eq:cases}
has $112{,}802$ elements. At every $(h,a)\in\mathcal C$, the recurrence
\eqref{eq:dp}--\eqref{eq:D} satisfies
\begin{equation}\label{eq:certified-maximum}
  D_h(a)\le
  \frac{184344982165038312381942179}{10^{26}}
  <\frac{36869}{20000}.
\end{equation}
\end{proposition}

\begin{proof}[Exact verification]
Run \verifier\ with the supplied \texttt{certificate.txt}. The program
first checks the nonnegativity and normalization of the bid weights and
sampling masses, the ordering of the multipliers and thresholds, and
$u_{120}/(u_{120}-1)<\rho$. It verifies all inequalities
\eqref{eq:cdf-checks} by exact integer comparisons.

For each $h$, it constructs $\paramset_h$ using precisely
\eqref{eq:endpoints}, discards values outside $I_h$, and removes duplicate
values within that set. For every resulting pair $(h,a)$, it constructs the
exact rational radius set \eqref{eq:radius-set} and evaluates an
outward-rounded integer version of \eqref{eq:dp}. As proved in
\cref{app:verification}, this gives an upper bound on $D_h(a)$; all
threshold comparisons use the exact radii, not their rounded values.
The computation uses no floating-point arithmetic.

There are $112{,}802$ endpoint cases. The largest certified upper bound is
the fraction in \eqref{eq:certified-maximum}, attained in this computation
at $h=84$, $a=869/405$, with selected path $869/405\longrightarrow0$.
The strict margin is
\begin{equation}\label{eq:margin}
  \frac{36869}{20000}
  -\frac{184344982165038312381942179}{10^{26}}
  =\frac{17834961687618057821}{10^{26}}>0.
\end{equation}
The full verification output is recorded in \texttt{verification.txt}.
The certificate data, integer recurrence, and reproduction commands are
specified in \cref{app:verification}.
\end{proof}

\begin{remark}[Meaning of the certified maximum]
The fraction in \eqref{eq:certified-maximum} is the largest
\emph{outward-rounded endpoint bound}. It is not asserted to be the exact
maximum of the unrounded dynamic programs or the optimum of the original
factor-revealing program. The theorem uses the declared target $\rho$,
which also determines the parameter domains $I_h$.
\end{remark}

\begin{proof}[Proof of \cref{thm:main}]
Start with any variables satisfying
\eqref{eq:row}--\eqref{eq:nonnegative}, without assuming normalization.
Replace the terminal radii by zero as in
\eqref{eq:terminal-replace}, and define $q_j$ and $C_j$ by
\eqref{eq:aj-qj} and \eqref{eq:credit}.

For a row with $d_j=0$, \cref{lem:zero-client} gives $q_j\le C_j$;
its cumulative-distribution hypothesis is supplied by
\cref{prop:certificate}.
For a row with $d_j>0$ and $a_j/d_j\le\rho$, \eqref{eq:q-a} gives
$q_j\le\rho d_j\le\rho d_j+C_j$.
For every other row, the reduction in \cref{sec:positive} assigns a pair
$(h,a)$ with $a\in I_h$ and gives \eqref{eq:row-excess-reduction}.
By \cref{lem:radius-reduction,lem:parameter-reduction,prop:certificate},
\[
  \frac{q_j-C_j}{d_j}
     \le D_h(a)
     \le\max_{a'\in\paramset_h}D_h(a')
     <\rho.
\]
Thus \eqref{eq:row-goal} holds for every row. Summing and applying
\eqref{eq:credit-budget} gives
\[
  \sum_{j=1}^n q_j
     \le\rho\sum_{j=1}^n d_j+\sum_{j=1}^n C_j
     \le\rho\sum_{j=1}^n d_j+f.
\]
The original transition-plus-terminal-radius expression is at most
$\sum_jq_j$, by the terminal-radius replacement. This proves
\eqref{eq:homogeneous-bound}.
Under \eqref{eq:normalization}, subtracting $f$ yields
$\Phi_\beta(r,f)\le\rho$. Taking the supremum over feasible variables and
then over $n$ proves \eqref{eq:main-bound}.
\end{proof}

\paragraph{Scope of the computational step.}
The verifier checks a fixed finite certificate. It does not optimize the
original program for small values of $n$, nor does it extrapolate from
numerical experiments. \Cref{lem:radius-reduction} covers every finite
radius path, and \cref{lem:parameter-reduction} covers the entire real
parameter domain, including its unbounded rays. These analytic results
are what make the finite computation sufficient for \cref{thm:main}.
This is a computer-assisted proof, not a proof-assistant formalization of
the analytic argument.

\begin{thebibliography}{1}
\bibitem{source-program}
Rares-Darius Buhai, Vincent Cohen-Addad, Henry Fleischmann, Anupam Gupta,
Antoni Jub\'es Monforte, Madhusudhan Reddy Pittu, Miltiadis Stouras,
Ola Svensson, and Lukas Vogl.
\newblock \emph{The Power of General Bids for $k$-Median}.
\newblock Supplied manuscript, Definitions~3.1 and~4.1.
\end{thebibliography}

\clearpage
\appendix
\section{Certificate data and exact verification}\label{app:verification}

All mathematical input to the computation is specified in this note:
the bid \eqref{eq:bid} and \cref{tab:bid}, the thresholds
\eqref{eq:times}, the masses in \cref{tab:masses}, the endpoint sets
\eqref{eq:endpoints}, and the recurrence \eqref{eq:dp}.
The files \texttt{certificate.txt} and \certfile\ contain the same
constants in machine-readable form. The former is the input to the C++
verifier, and the latter is used by the supplementary Python checks.

\begin{table}[H]
\centering
\caption{The rational sampling distribution. At threshold
$u_\ell=1+\ell/96$, the probability mass is $c_\ell/10^9$.
All $121$ integers, including zero masses, are listed.}
\label{tab:masses}
\small
\renewcommand{\arraystretch}{1.08}
\setlength{\tabcolsep}{7pt}
\begin{tabular}{rr@{\qquad}rr@{\qquad}rr@{\qquad}rr}
\toprule
$\ell$ & $c_\ell$ & $\ell$ & $c_\ell$ & $\ell$ & $c_\ell$ & $\ell$ & $c_\ell$\\
\midrule
0 & $24{,}087{,}553$ & 31 & $15{,}347{,}486$ & 62 & $5{,}430{,}994$ & 93 & $0$\\
1 & $23{,}567{,}894$ & 32 & $19{,}903{,}937$ & 63 & $5{,}394{,}373$ & 94 & $15{,}425{,}764$\\
2 & $23{,}315{,}363$ & 33 & $10{,}306{,}194$ & 64 & $5{,}276{,}652$ & 95 & $0$\\
3 & $22{,}706{,}147$ & 34 & $14{,}496{,}475$ & 65 & $5{,}164{,}893$ & 96 & $0$\\
4 & $22{,}224{,}377$ & 35 & $13{,}786{,}530$ & 66 & $5{,}037{,}775$ & 97 & $0$\\
5 & $21{,}750{,}480$ & 36 & $13{,}032{,}292$ & 67 & $4{,}928{,}361$ & 98 & $0$\\
6 & $21{,}230{,}145$ & 37 & $12{,}186{,}845$ & 68 & $4{,}811{,}734$ & 99 & $0$\\
7 & $20{,}726{,}429$ & 38 & $11{,}086{,}606$ & 69 & $4{,}716{,}895$ & 100 & $0$\\
8 & $20{,}611{,}532$ & 39 & $10{,}315{,}351$ & 70 & $4{,}616{,}748$ & 101 & $0$\\
9 & $20{,}155{,}272$ & 40 & $10{,}140{,}726$ & 71 & $4{,}537{,}378$ & 102 & $3{,}302{,}130$\\
10 & $19{,}733{,}910$ & 41 & $9{,}982{,}272$ & 72 & $4{,}442{,}037$ & 103 & $0$\\
11 & $19{,}616{,}405$ & 42 & $9{,}828{,}489$ & 73 & $4{,}371{,}224$ & 104 & $0$\\
12 & $19{,}195{,}203$ & 43 & $9{,}679{,}177$ & 74 & $4{,}289{,}205$ & 105 & $0$\\
13 & $18{,}952{,}661$ & 44 & $9{,}534{,}145$ & 75 & $4{,}230{,}208$ & 106 & $0$\\
14 & $20{,}108{,}443$ & 45 & $8{,}588{,}766$ & 76 & $4{,}158{,}464$ & 107 & $0$\\
15 & $18{,}554{,}301$ & 46 & $5{,}363{,}641$ & 77 & $4{,}101{,}646$ & 108 & $0$\\
16 & $16{,}736{,}136$ & 47 & $5{,}440{,}727$ & 78 & $4{,}036{,}047$ & 109 & $0$\\
17 & $17{,}512{,}178$ & 48 & $5{,}508{,}246$ & 79 & $3{,}984{,}451$ & 110 & $0$\\
18 & $17{,}533{,}687$ & 49 & $5{,}559{,}267$ & 80 & $3{,}920{,}032$ & 111 & $0$\\
19 & $17{,}195{,}813$ & 50 & $5{,}592{,}099$ & 81 & $3{,}878{,}363$ & 112 & $0$\\
20 & $17{,}124{,}860$ & 51 & $5{,}609{,}257$ & 82 & $3{,}822{,}612$ & 113 & $0$\\
21 & $16{,}945{,}175$ & 52 & $5{,}614{,}263$ & 83 & $13{,}251{,}875$ & 114 & $0$\\
22 & $16{,}637{,}025$ & 53 & $5{,}615{,}387$ & 84 & $0$ & 115 & $0$\\
23 & $16{,}595{,}629$ & 54 & $5{,}615{,}841$ & 85 & $0$ & 116 & $0$\\
24 & $16{,}390{,}990$ & 55 & $5{,}609{,}762$ & 86 & $8{,}798{,}431$ & 117 & $0$\\
25 & $16{,}312{,}976$ & 56 & $5{,}603{,}100$ & 87 & $0$ & 118 & $0$\\
26 & $16{,}176{,}735$ & 57 & $5{,}600{,}838$ & 88 & $0$ & 119 & $0$\\
27 & $15{,}998{,}131$ & 58 & $5{,}591{,}664$ & 89 & $8{,}086{,}571$ & 120 & $0$\\
28 & $15{,}833{,}420$ & 59 & $5{,}570{,}605$ & 90 & $0$ &  & \\
29 & $15{,}633{,}516$ & 60 & $5{,}519{,}377$ & 91 & $0$ &  & \\
30 & $15{,}399{,}233$ & 61 & $5{,}475{,}884$ & 92 & $4{,}318{,}269$ &  & \\
\bottomrule
\end{tabular}
\end{table}

\Needspace{10\baselineskip}
\subsection{The integer recurrence}

Let
\begin{equation}\label{eq:integer-constants}
  B=10^8,\qquad D=10^9,\qquad R=10^9,\qquad G=4500,
\end{equation}
so that $\gamma_s=p_s/G$ with each $p_s$ an integer. For nonnegative
integers $X,Y$, define
\begin{equation}\label{eq:integer-bid}
  N(X,Y)=\sum_{s=1}^{13}W_s\max\{GX-p_sY,0\}.
\end{equation}
Then
\begin{equation}\label{eq:integer-bid-identity}
  \beta(X/R,Y/R)=\frac{N(X,Y)}{BGR}.
\end{equation}
The implementation evaluates \eqref{eq:integer-bid} using prefix sums of
$W_s$ and $W_sp_s$. Since the $p_s$ are increasing, the positive terms
form a prefix; an exact integer comparison locates its endpoint.

Fix an endpoint case $(h,a)$ and retain its exact rational states $v_k$.
For every nonnegative rational $v$, put
\[
  v^- =\lfloor Rv\rfloor,\qquad v^+=\lceil Rv\rceil.
\]
Define the integer threshold masses
\begin{equation}\label{eq:integer-masses}
  s_k=\sum_{\substack{\ell\ge h\\b_\ell(a)>v_k}}c_\ell
      =D S_h(v_k;a).
\end{equation}
In particular, $s_p-s_k\ge0$ whenever $p<k$.
The integer recurrence is
\begin{align}
  \widehat V_0&=0,\notag\\
  \widehat V_k
    &=\max_{p<k}\left\{
       D N(v_k^+,v_p^-)
       -(s_p-s_k)N(v_k^-,R)
       +\widehat V_p
       \right\},\label{eq:integer-dp}\\
  \widehat A_h(a)
    &=\sum_{\ell<h}c_\ell
       N\bigl(\lfloor R M_h(a)u_\ell\rfloor,R\bigr).
       \label{eq:integer-active}
\end{align}
Its output, scaled back to the original units, is
\begin{equation}\label{eq:rounded-D}
  \widehat D_h(a)
   =\frac{\widehat V_K-\widehat A_h(a)}{BGRD}.
\end{equation}

\begin{lemma}[Directed-rounding guarantee]\label[lemma]{lem:directed-rounding}
For every endpoint case $(h,a)$,
\begin{equation}\label{eq:integer-upper}
  D_h(a)\le\widehat D_h(a).
\end{equation}
\end{lemma}

\begin{proof}
Monotonicity in the two arguments gives
\[
  \beta(v_k,v_p)
       \le\beta(v_k^+/R,v_p^-/R)
       =\frac{N(v_k^+,v_p^-)}{BGR},
\]
whereas
\[
  \beta(v_k,1)\ge\frac{N(v_k^-,R)}{BGR}.
\]
Since the mass $s_p-s_k$ is nonnegative, these are upper bounds on the
positive transition payment and lower bounds on the reference bid being
subtracted. Induction on $k$ in \eqref{eq:dp} and
\eqref{eq:integer-dp} yields $BGRD\,V_k\le\widehat V_k$.
Similarly, \eqref{eq:integer-active} satisfies
$\widehat A_h(a)\le BGRD\,A_h(a)$, since every active reference argument is
rounded downward. Subtracting proves \eqref{eq:integer-upper}.

The comparisons in \eqref{eq:integer-masses} use the exact rational
values $b_\ell(a)$ and $v_k$. They are not affected by rounding, so the
strict-crossing convention is preserved exactly.
\end{proof}

This calculation rounds in a specified direction using integers; it is
not a floating-point approximation with an estimated error. The final
comparison with $\rho$ is made after multiplication by the positive exact
denominators.

\subsection{Input files, commands, and verification record}

The header of \texttt{certificate.txt} supplies the numerator and denominator
of $\rho$, then $B,D,R$, then the numbers of hinges and thresholds.
Its exact contents are
\begin{verbatim}
36869 20000
100000000 1000000000 1000000000
13 121
\end{verbatim}
Line breaks serve only as whitespace in the input.
They are followed by $13$ lines containing the numerator and denominator
of $\gamma_s$ and the integer $W_s$, then $121$ lines containing the
numerator and denominator of $u_\ell$ and the integer $c_\ell$.
Thus the text certificate can be reconstructed entirely from this note.
The verifier computes $G$ as the least common multiple of the hinge
denominators.

Place \verifier\ and \texttt{certificate.txt} in the same directory,
and use a C++17 compiler supporting \texttt{\_\_int128}, with the
Boost.Multiprecision headers available:
\begin{verbatim}
g++ -O3 -std=c++17 verify.cpp -o verify
./verify certificate.txt
\end{verbatim}
The successful output is
\begin{verbatim}
PASS: exact all-dimensions threshold certificate
Target: 36869/20000
Mixture components: 13
Sampling thresholds: 121
Exact CDF inequalities: 120
Endpoint cases: 112802
Largest outward-rounded upper bound:
  184344982165038312381942179/100000000000000000000000000
Margin:
  17834961687618057821/100000000000000000000000000
Worst checked case: h=84, a=869/405
Worst checked path: 869/405 0
\end{verbatim}
Here the two longest lines have been broken for typesetting only.
The function \texttt{endpoints} implements \eqref{eq:endpoints}, and
\texttt{check\_case} implements \eqref{eq:integer-dp}--\eqref{eq:rounded-D}.

The program uses reduced rational numbers for endpoints and radius states,
signed $128$-bit integers for the recurrence, and arbitrary-precision
integers for cumulative-mass and final target comparisons. It checks that
rational components fit their storage types and uses an
arbitrary-precision magnitude bound before evaluating each recurrence.
A certificate outside the checked arithmetic range is rejected.

Two supplementary scripts provide independent rational checks:
\begin{verbatim}
python cross_check.py
python verify_lower_bound.py
\end{verbatim}
The first script reads \certfile, checks the cumulative-mass conditions
using Python's \texttt{Fraction} arithmetic, and compares selected
\emph{unrounded} dynamic programs with the C++ upper bounds. It also
checks the one-outlier diagnostic discussed in
\cref{app:lower}. These selected checks are supplementary; the C++
verification, not the sample of Python cases, checks all $112{,}802$
endpoints. The second script checks the exact lower-bound values.
Their outputs are supplied in \texttt{independent\_checks.txt} and
\texttt{lower\_bound\_checks.txt}.

\section{A nearly matching feasible family}\label{app:lower}

This appendix records a lower bound for the same optimization problem.
It is not used in the upper-bound proof.

\begin{proposition}[One-outlier lower bound]\label[proposition]{prop:lower}
For the bid \eqref{eq:bid},
\begin{equation}\label{eq:lower}
  \sup_{n\ge1}\FR_n(\beta)
     \ge\frac{8295524183431}{4500000000000}
     =1.843449818540222\ldots.
\end{equation}
\end{proposition}

\begin{proof}
More generally, for any $z>0$ define
\[
  L_\beta(z)=z+\beta(z+1,z)-\beta(z,1).
\]
We construct feasible points whose objectives converge to $L_\beta(z)$.
Let $M$ be a positive integer with $M\ge z$, and take $n=M+1$ clients.
Set
\[
  d_1=1,\qquad d_2=\cdots=d_n=0,\qquad
  T=z+1,\qquad a=T-\frac zM,\qquad
  f=MT+\beta(z,1).
\]
Define the radii by
\begin{align*}
  r_{1,1}&=a,&
  r_{1,i}&=z\quad(2\le i\le n),&
  r_{1,n+1}&=0,\\
  r_{j,i}&=T\quad(2\le j\le i\le n),&&
  r_{j,n+1}&=0\quad(2\le j\le n).
\end{align*}
Normalization and nonnegativity hold. Since $M\ge z$, we have $a\ge z$,
so every row is nonincreasing. A triangle constraint between the first row
and a later diagonal is tight, because $T=z+1$; a triangle constraint
between two zero-distance rows is also tight. The only nonzero difference
between two earlier rows in a nonterminal column is $T-z=1$, exactly the
sum of their reference distances. All terminal radii are zero. Thus all
triangle and metric constraints hold.

For every facility constraint at $i\ge2$, the left-hand side is
$\beta(z,1)+MT=f$. At $i=1$, budget feasibility and $a\le T$ give
\[
  \beta(a,1)+Ma
    \le z+M\left(T-\frac zM\right)
    =MT\le f.
\]
Hence the point is feasible for the full program. Its objective is
\[
  z+\beta\left(z+1-\frac zM,z\right)-\beta(z,1).
\]
By continuity, this tends to $L_\beta(z)$ as $M\to\infty$.

For the bid in \cref{tab:bid}, take $z=17/10$. Direct rational evaluation
of \eqref{eq:bid} gives
\[
  L_\beta(17/10)
    =\frac{8295524183431}{4500000000000},
\]
proving \eqref{eq:lower}. These evaluations are checked by
\texttt{verify\_lower\_bound.py}.
\end{proof}

\end{document}
